\documentclass[a4paper,10pt]{article}
\usepackage[spanish]{babel}
\usepackage{
			ucs
			,amsfonts
			,amssymb
			,acronym
			,verbatim
			,rotating
			,graphicx
			,enumerate
			,fancyhdr
            ,mathrsfs
            ,array
			}

\usepackage{tikz}
\usetikzlibrary{svg.path}

\pagestyle{fancy}
\addtolength{\textheight}{1.0in}
\addtolength{\textwidth}{0.5in}
\addtolength{\voffset}{-1.0in}

\fancyhf{}
\fancyhead[L]{{\bfseries Programaci\'on I Curso 2012. Pr\'actica 5 \hfill Funciones y Procedimientos}}
\fancyfoot[LE,RO]{\thepage}

\fancypagestyle{firstpage}{
\fancyhf{}

\fancyhead[L]{\includegraphics[width=1.75cm]{ucalp.jpg}  \\
\small{UCALP. Facultad de Ciencias Exactas e Ingenier\'ia} \\
{\bfseries Programaci\'on I Curso 2012. Pr\'actica 5 \hfill Funciones y Procedimientos}}
\fancyfoot[LE,RO]{\thepage}
}

\thispagestyle{firstpage}

\begin{document}

~\\ \\ \\ \\

\begin{enumerate}

\item (Pr\'actica 3) Desarrollar un programa para, dado un conjunto de n\'umeros enteros, escribir aquellos que sean divisibles por $I$. Crear una funci\'on que verifique si un n\'umero es divisible por $I$.   
    
\item (Pr\'actica 3) Desarrolle una funci\'on/procedimiento que permita saber si un n\'umero ingresado por teclado es par (positivo y negativo) o impar (positivo y negativo).

\item (Pr\'actica 4) Realizar un programa que permita almacenar un listado de palabras en un vector y luego imprimirlas en pantalla. El programa debe contener al menos 2 funciones/procedimientos: uno para cargar los datos y otra para mostrarlos.

\item (Pr\'actica 4) Se desea realizar una funci\'on/procedimiento que nos permita buscar una palabra dentro de una frase y cuantas veces se repite.

\item (Pr\'actica 4) Escriba un par de funciones que codifique una contrase\~na de 8 d\'igitos. El algoritmo deber\'a sustituir cada vocal por un n\'umero. $(A=1, E=2, O=3, I=4, U=5)$. Se desea realizar el codificador y decodificador.


\item (Pr\'actica 4) Dados un vector , una matriz y un valor, escribir funciones/procedimientos que puedan:
    \begin{itemize}
        \item[a-] Devolver la cantidad de ocurrencias de dicho valor en el vector.
        \item[b-] Devolver la cantidad de ocurrencias de dicho valor en la matriz.
        \item[c-] Devolver el mismo vector donde todas las ocurrencias de dicho valor han sido suprimidas.
        \item[c-] Devolver todas las ocurrencias del vector dentro de la matriz.
    \end{itemize}

\item (Pr\'actica Integradora) Realizar una funci\'on que devuelva el resultado de la interseccio\'n de dos conjuntos.

\item (Pr\'actica Integradora) Realizar una funci\'on/procediemnto que multiplique una matriz por un escalar.

\item (Pr\'actica Integradora) Realizar una funci\'on/procedimiento que retorne el promedio de un vector de n\'umeros.

\end{enumerate}
\end{document}
